\chapter{Use Cases}
\label{chap-use-case}

\phantomsection
\fancyhead[LE]{\thepage}
\fancyhead[RE]{\textbf{Hydrological Twin - Architectural Decomposition}}
\fancyhead[RO]{\thepage}
\fancyhead[LO]{\textbf{Hydrological Twin - Architectural Decomposition}}
\renewcommand{\headrulewidth}{1pt}

\setlength{\parindent}{0cm}
\setlength{\parskip}{9pt}

% ---------------------------------------------------------------------

\section{Scientific Integrity of Simulation}

Instanciating a HydroTwin requires to check its integrity. It requires private-key level


HydroTwin performs the certification of its integrity, meaning it can reproduce the simulation (multi-model outputs) based on the provided information. 


\section{Simulation rendering and basic statistics - CaWaQS-ViZ}

public

Reexpliquer ce qui est dit précédemment mais en s'inspirant de la notice cawaqs-viz

\section{Generate a submodel}

Extraction of a subTwin -> metadata, subModels that run, Simulation OUTPUTS. private-key level - nested functionality. HydroTwin must contain a pointeur towards its subTwin. Do be done via the global registry, or the subTwin is contained in its mother class. I prefer option2

Cas d'usage typique un territoire en son bassin, un champ captant en son territoire

\section{Analysis}

private-key 

parameter self-fitting
uncertainty analysis

based on bayesian inversion, deep learning, generating